% Generated from paper/full_draft. Edit the Markdown source or migration script.

Knowledge-augmented agents must select enough evidence to support answer quality while limiting latency, noise, and final context cost. This trade-off is especially difficult for vertical-domain data, where relevance is shaped by domain terminology, local semantic neighborhoods, workflow structure, and evolving user intent. We propose IntentWeight, a feedback-guided evidence selection controller motivated by a piecewise relevance-manifold assumption. Here, the assumption means that query-document relevance in vertical domains often exposes exploitable local structure rather than a uniform embedding-space distribution. IntentWeight combines dense semantic retrieval, BM25 lexical recall, and cluster-local routing, and uses trust-weighted LinUCB to learn route preferences from simulated feedback. A confidence-based final context policy then compacts the selected evidence sent to the generator while preserving dense fallback under low confidence. We instantiate this framework in retrieval-augmented question answering on LoTTE technology/search from 100k to 638k corpus chunks. Frozen calibration/test budget policies save 6-18\% final LLM evidence-context input tokens while outperforming dense-only adaptive truncation in $\mathrm{Hit@10}$, with scale-dependent non-inferiority. A conservative confidence policy provides a stable 4.7-5.3\% saving baseline. Results generalize to a second LoTTE domain with domain-calibrated compression, and simulated feedback recovers a meaningful fraction of compression-induced tail failures. IntentWeight is therefore not a universal dense replacement, but an adaptive quality-cost and recovery controller for structured domain evidence.
